% Print styling for the prayer blocks. The web book's look reduced to what
% paper can carry: an indented block, the text at reading size, the
% transliteration under it in smaller grey italic.

\newlength{\prayerindent}
\setlength{\prayerindent}{1.2em}

% \interlinepenalty keeps a prayer whole: a text and its transliteration are
% one unit, and a version split across a page turn is hard to read against
% the others. LaTeX moves the whole block down instead.
\newenvironment{prayer}{%
  \par\medskip
  \begingroup
  \setlength{\parindent}{0pt}%
  \setlength{\leftskip}{\prayerindent}%
  \setlength{\parskip}{0pt}%
  \interlinepenalty=10000
  \predisplaypenalty=10000
}{%
  \par
  \endgroup
  \medskip
}

% Transliteration: quieter than the text it glosses, and never hyphenated,
% so a romanised Greek line breaks where the Greek does.
\newcommand{\prayertr}[1]{%
  {\small\itshape\textcolor{black!55}{\hyphenpenalty=10000 #1}}%
}

% KOMA sets headings in sans by default, which leaves the book in two
% unrelated faces. Everything here is EB Garamond, as on the web.
\addtokomafont{disposition}{\rmfamily}
\addtokomafont{chapter}{\rmfamily\bfseries}

% A5 pages cannot spare the drop that scrbook gives a chapter opening.
\RedeclareSectionCommand[beforeskip=0pt, afterskip=1.2\baselineskip]{chapter}
\RedeclareSectionCommand[beforeskip=-1.4\baselineskip, afterskip=.5\baselineskip]{section}
